% Ch2_7, slide 5: two conductors carrying +Q and -Q, the field lines between them and the source V12
\begin{tikzpicture}[scale=1]
  \coordinate (A) at (-1.6,0.2); \coordinate (B) at (1.7,0.3);
  \fill[ELfillB] plot[smooth cycle,tension=0.8] coordinates{(-2.4,-0.3)(-2.0,0.9)(-1.4,1.3)(-0.9,0.9)(-0.8,0.1)(-1.1,-0.6)(-1.8,-0.8)};
  \draw[ELcField!80,line width=0.7pt] plot[smooth cycle,tension=0.8] coordinates{(-2.4,-0.3)(-2.0,0.9)(-1.4,1.3)(-0.9,0.9)(-0.8,0.1)(-1.1,-0.6)(-1.8,-0.8)};
  \fill[ELfillB] plot[smooth cycle,tension=0.8] coordinates{(1.2,0.3)(1.5,0.9)(1.95,0.85)(2.2,0.3)(1.95,-0.25)(1.5,-0.3)};
  \draw[ELcField!80,line width=0.7pt] plot[smooth cycle,tension=0.8] coordinates{(1.2,0.3)(1.5,0.9)(1.95,0.85)(2.2,0.3)(1.95,-0.25)(1.5,-0.3)};
  \node[ELlabs] at (-1.6,0.2) {$+Q$}; \node[ELlabs] at (1.7,0.3) {$-Q$};
  \foreach \y/\b in {0.95/30,0.55/12,0.2/0,-0.15/-10,-0.5/-25}{\draw[ELcSource,dashed,line width=0.6pt,ELmid] (-0.85,\y) to[out=\b,in=180-\b] (1.2,{0.3+0.45*\y});}
  \draw[ELcSource,dashed,line width=0.6pt,ELmid] (-1.4,1.3) to[out=70,in=110] (1.6,0.9);
  \foreach \a/\x/\y in {150/-2.2/0.7,200/-2.35/-0.4,250/-1.6/-0.75}{\draw[ELcSource,dashed,line width=0.6pt,-{Stealth[length=4pt]}] (\x,\y)--++(\a:0.6);}
  \foreach \a/\x/\y in {20/2.15/0.55,-20/2.1/0.0}{\draw[ELcSource,dashed,line width=0.6pt,{Stealth[length=4pt]}-] (\x,\y)--++(\a:0.55);}
  \draw[ELink,line width=0.7pt] (-2.4,-0.3)--(-2.9,-0.3) (2.2,0.3)--(2.9,0.3)--(2.9,-1.7)--(0.35,-1.7) (-0.15,-1.7)--(-2.9,-1.7)--(-2.9,-0.3);
  \draw[ELink,line width=0.9pt] (-0.15,-1.95)--(-0.15,-1.45) (0.05,-1.85)--(0.05,-1.55) (0.2,-1.95)--(0.2,-1.45) (0.35,-1.85)--(0.35,-1.55);
  \node[ELlabs,anchor=north] at (0.1,-1.98) {$V_{12}$};
  \node[ELlabs,anchor=south] at (-0.2,-1.5) {$+$}; \node[ELlabs,anchor=south] at (0.45,-1.5) {$-$};
\end{tikzpicture}
